\documentclass[11pt]{article}

% ======================
% Packages
% ======================
\usepackage[utf8]{inputenc}
\usepackage[T1]{fontenc}
\usepackage{lmodern}
\usepackage{amsmath, amssymb, amsthm}
\usepackage{geometry}
\usepackage{listings}
\usepackage{xcolor}
\usepackage{graphicx}
\usepackage{algorithm}
\usepackage{algpseudocode}
\usepackage{hyperref}
\usepackage{enumitem}
\usepackage{enumerate}

% ======================
% Page Layout
% ======================
\geometry{a4paper, margin=1in}
\setlength{\parindent}{0pt}
\setlength{\parskip}{1em}

% ======================
% Colors & Code Styling
% ======================
\definecolor{codegreen}{rgb}{0,0.6,0}
\definecolor{codegray}{rgb}{0.5,0.5,0.5}
\definecolor{codepurple}{rgb}{0.58,0,0.82}
\definecolor{backcolour}{rgb}{0.95,0.95,0.92}

\lstdefinestyle{mystyle}{
    backgroundcolor=\color{backcolour},
    commentstyle=\color{codegreen},
    keywordstyle=\color{magenta},
    numberstyle=\tiny\color{codegray},
    stringstyle=\color{codepurple},
    basicstyle=\ttfamily\small,
    breakatwhitespace=false,
    breaklines=true,
    captionpos=b,
    keepspaces=true,
    numbers=left,
    numbersep=5pt,
    showspaces=false,
    showstringspaces=false,
    showtabs=false,
    tabsize=2
}

\lstset{style=mystyle}

% ======================
% Theorem-like Environments
% ======================
\newtheorem{problem}{Problem}
\newtheorem{subproblem}{Subproblem}[problem]
\newtheorem{solution}{Solution}[problem]

% ======================
% Title Info
% ======================
\title{\textbf{Fundamental Algorithm Techniques} \\ Problem Set \#6} % \author{}
\date{\textcolor{red}{Review on Mars 12}} %Due: November 01, 2025}

% ======================
% Document
% ======================
\begin{document}

\maketitle



\begin{problem}[Facebook Interview, 10/10 pts]
We consider trees of $n$ children,  each with weight $1/n$ of their parent's weight.
\begin{enumerate}
    \item Write the general class for this object 
    \item Generate a tree of depth $N=3$, with initial parent tree of weight $1/n$
    \item create a depth first recursive function visiting each node and summing up the weights. Make sure it returns 1 for various $n$'s! 
    \item Same with breadth first, check also 1!
    \item Same as above for both searches, but each time you reach a node, flip the value sign. Make sure you get 1 and -1 after both first and second searches (run a test with fixed $n$)!
    \item Write recursive and non recursive versions of breadth first.
    \item Whilst depth first can be implemented recursively, it is not recommended for breadth first! Explain why? 
    \item Post discussion and code to github and in time!
\end{enumerate}

\end{problem}




\end{document}

